
\medskip
\noindent\SGAInlineRefTarget{sga3:iii:paragraph:0-1-7}\textbf{Reminders 0.1.7.}\footnote{Editor's note: This
paragraph has been added (see also [DG70], \S I.4, nos.~1--2).}
In this paragraph, given an \(S\)-scheme \(X\), we ``recall'' certain
functorial properties of the module of differentials
\(\Omega^1_{X/S}\) and of the first infinitesimal neighborhood of the
diagonal, \(\Delta^{(1)}_{X/S}\). These properties follow readily from
EGA~IV\textsubscript{4}, \S\S16.1--16.4, but do not appear there
explicitly.

\smallskip
\noindent\textit{(a)} Let us begin by recalling the following facts
(cf. EGA~II, \S\S1.2--1.5). Let \(g:Y\to X\) be a morphism of schemes,
let \(\pi:X'\to X\) be an affine \(X\)-scheme, and let
\(\mathscr B=\pi_*(\mathscr O_{X'})\) be the corresponding
quasi-coherent \(\mathscr O_X\)-algebra. Then the \(Y\)-scheme
\(Y\times_X X'\) is affine and corresponds to the quasi-coherent
\(\mathscr O_Y\)-algebra \(g^*(\mathscr B)\), and there is a
commutative diagram of bijections:

\SGAThreeNativeDiagram[0.82\textwidth]
  {native_diagrams/exp3/diagram_003_exp3_p005_reminders017_affine_bijections.tex}
  {The affine-scheme/algebra bijections after base change by \(g\).}
  {fig:exp3-reminders-017-affine-bijections}
  {Polo--Gille Exp3-14oct24.pdf, component-local p. 5; combined-reader p. 115}

Moreover, these bijections are functorial in the pair \((X,X')\). More
precisely, let \(W'\) be an affine \(W\)-scheme corresponding to the
quasi-coherent \(\mathscr O_W\)-algebra \(\mathscr A\), and suppose that
one has a commutative diagram of morphisms of schemes as follows.

\SGAThreeNativeDiagram[0.38\textwidth]
  {native_diagrams/exp3/diagram_004_exp3_p005_reminders017_scheme_functoriality.tex}
  {The scheme morphisms used to express functoriality in \((X,X')\).}
  {fig:exp3-reminders-017-scheme-functoriality}
  {Polo--Gille Exp3-14oct24.pdf, component-local p. 5; combined-reader p. 115}

Denote by
\[
 \varphi:\mathscr A\longrightarrow f_*(\mathscr B),
 \qquad
 \varphi^\sharp:f^*(\mathscr A)\longrightarrow\mathscr B
\]
the algebra morphisms associated with \(f'\), and by
\[
 \theta:\mathscr B\longrightarrow g_*(\mathscr O_Y),
 \qquad
 \theta^\sharp:g^*(\mathscr B)\longrightarrow\mathscr O_Y
\]
those associated with a variable \(X\)-morphism \(g':Y\to X'\).
Then the following diagram is commutative, where \(Y\) is viewed as a
\(W\)-scheme through \(f\circ g\).

\SGAThreeNativeDiagram[0.96\textwidth]
  {native_diagrams/exp3/diagram_005_exp3_p006_reminders017_hom_functoriality.tex}
  {Functoriality of the affine-scheme/algebra correspondences.}
  {fig:exp3-reminders-017-hom-functoriality}
  {Polo--Gille Exp3-14oct24.pdf, component-local p. 6; combined-reader p. 116}

\smallskip
\noindent\SGAInlineRefTarget{sga3:iii:paragraph:0-1-7:item:b}\textit{(b)} Let now \(X\) be an \(S\)-scheme. Let
\(\Omega^1_{X/S}\) be the module of differentials of \(X\) over \(S\),
and let \(\Delta^{(1)}_{X/S}\) be the first infinitesimal neighborhood
of the diagonal immersion
\(\delta_X:X\to X\times_S X\). It is a subscheme of
\(X\times_S X\), in which the diagonal \(\Delta_{X/S}\) is a closed
subscheme. Write \(\operatorname{pr}^i_X\), \(i=1,2\), for the two
projections \(X\times_S X\to X\), and write \(\pi_X\) for the
restriction of \(\operatorname{pr}^1_X\) to
\(\Delta^{(1)}_{X/S}\).

On the one hand, every morphism \(f:X\to W\) of \(S\)-schemes induces
an \(S\)-morphism
\[
 \Delta^{(1)}f:\Delta^{(1)}_{X/S}\longrightarrow
 \Delta^{(1)}_{W/S}
\]
for which the following diagram is commutative.

\SGAThreeNativeDiagram[0.86\textwidth]
  {native_diagrams/exp3/diagram_006_exp3_p006_reminders017_delta_functoriality.tex}
  {Functoriality of the first infinitesimal neighborhood of the diagonal.}
  {fig:exp3-reminders-017-delta-functoriality}
  {Polo--Gille Exp3-14oct24.pdf, component-local p. 6; combined-reader p. 116}

On the other hand, through the projection \(\pi_X\), the scheme
\(\Delta^{(1)}_{X/S}\) is an affine \(X\)-scheme: it is the spectrum
of the augmented quasi-coherent \(\mathscr O_X\)-algebra
\[
 \mathscr P^1_{X/S}=\mathscr O_X\oplus\Omega^1_{X/S},
\]
where \(\Omega^1_{X/S}\) is a square-zero ideal. The augmentation is
the \(\mathscr O_X\)-algebra morphism
\[
 \varepsilon_X:\mathscr P^1_{X/S}\longrightarrow\mathscr O_X
\]
that vanishes on \(\Omega^1_{X/S}\) and corresponds to the closed
immersion \(\delta_X:X\hookrightarrow\Delta^{(1)}_{X/S}\). Hence every
morphism of \(S\)-schemes \(f:X\to W\) induces a morphism of augmented
\(\mathscr O_X\)-algebras
\[
 f^*(\mathscr P^1_{W/S})
 =\mathscr O_X\oplus f^*(\Omega^1_{W/S})
 \longrightarrow
 \mathscr P^1_{X/S}=\mathscr O_X\oplus\Omega^1_{X/S},
\]
or, equivalently, a morphism of \(\mathscr O_X\)-modules
\[
 f_{X/W/S}:f^*(\Omega^1_{W/S})\longrightarrow\Omega^1_{X/S};
\]
see EGA~IV\textsubscript{4}, (16.4.3.6), and (16.4.18.2) for the
notation \(f_{X/W/S}\).

Since \(\pi_X:\Delta^{(1)}_{X/S}\to X\) is affine, part~(a) shows
that \(\Delta^{(1)}f\) is completely determined by \(f_{X/W/S}\).
For every \(X\)-scheme \(g:Y\to X\), the set
\[
 \operatorname{Hom}_X(Y,\Delta^{(1)}_{X/S})
 \simeq
 \operatorname{Hom}_{\mathscr O_Y\text{-}\mathrm{alg}}
 \bigl(\mathscr O_Y\oplus g^*(\Omega^1_{X/S}),\mathscr O_Y\bigr)
\]
identifies with a subset of
\(\operatorname{Hom}_{\mathscr O_Y}
 (g^*(\Omega^1_{X/S}),\mathscr O_Y)\), namely
\[
 \widetilde{\operatorname{Hom}}_{\mathscr O_Y}
 \bigl(g^*(\Omega^1_{X/S}),\mathscr O_Y\bigr),
\]
the subset formed by the \(\mathscr O_Y\)-morphisms
\(\psi:g^*(\Omega^1_{X/S})\to\mathscr O_Y\) for which
\(\operatorname{Im}(\psi)\) is a square-zero ideal of
\(\mathscr O_Y\).\footnote{Editor's note: In the proof of SGA~1,
III~5.1, the sentence ``But the algebra homomorphisms
\(\mathscr A\to\mathscr O_Y\) correspond \ldots'' must be corrected
accordingly; the remainder of the cited passage is unchanged.}

Consequently, apply part~(a) to the following diagram.

\SGAThreeNativeDiagram[0.48\textwidth]
  {native_diagrams/exp3/diagram_007_exp3_p007_reminders017_delta_application.tex}
  {The diagram to which the affine functoriality statement is applied.}
  {fig:exp3-reminders-017-delta-application}
  {Polo--Gille Exp3-14oct24.pdf, component-local p. 7; combined-reader p. 117}

Taking into account that \(\Delta^{(1)}f\) is the restriction of
\(f\times f\) to \(\Delta^{(1)}_{X/S}\), one obtains the following
commutative diagram, functorial in the \(X\)-scheme \(g:Y\to X\).

\SGAThreeNativeDiagram[0.97\textwidth]
  {native_diagrams/exp3/diagram_008_exp3_p007_reminders017_star_diagram.tex}
  {The functorial comparison of first-neighborhood morphisms and differential maps, equation \(0.1.7\,(*)\).}
  {fig:exp3-reminders-017-star}
  {Polo--Gille Exp3-14oct24.pdf, component-local p. 7; combined-reader p. 117}

\medskip
\noindent\SGAInlineRefTarget{sga3:iii:remark:0-1-7-1}\textbf{Remark 0.1.7.1.}
Let us conclude this paragraph with the following remark, which will be
useful later (cf. \SGARefLink{sga3:iii:corollary:0-1-10}{0.1.10}). Denote by \(\widetilde L_X\) the
\(X\)-functor that associates to every \(X\)-scheme \(g:Y\to X\) the
set
\[
 \operatorname{Hom}_{\mathscr O_Y}
 \bigl(g^*(\Omega^1_{X/S}),\mathscr O_Y\bigr),
\]
and denote by \(\widetilde L_f:\widetilde L_X\to\widetilde L_W\) the
morphism of functors defined above: it sends
\(\psi\in\widetilde L_X(Y)\) to
\(\psi\circ g^*(f_{X/W/S})\). The preceding discussion gives a
commutative diagram of functors.

\SGAThreeNativeDiagram[0.74\textwidth]
  {native_diagrams/exp3/diagram_009_exp3_p007_remark0171_functor_diagram.tex}
  {Functoriality of \(\widetilde L_X\) and the first infinitesimal neighborhood.}
  {fig:exp3-remark-0171-functor-diagram}
  {Polo--Gille Exp3-14oct24.pdf, component-local p. 7; combined-reader p. 117}

\medskip
\noindent\SGAInlineRefTarget{sga3:iii:theorem:0-1-8}\textbf{Theorem 0.1.8.} (SGA~1, III~5.1)\footnote{Editor's
note: The statement and proof of SGA~1, III~5.1 are reproduced here,
since they are needed below to prove \SGARefLink{sga3:iii:proposition:0-2}{Proposition~0.2}; see already the
following \SGARefLink{sga3:iii:corollary:0-1-9}{Corollary~0.1.9}. Moreover, \(\mathscr J\) (respectively
\(\mathscr I\)) is another calligraphic form of the letter \(J\)
(respectively \(I\)). Returning to the earlier notation, with
\(\mathscr J\subset\mathscr I\) ideals of \(\mathscr O_S\) such that
\(\mathscr I\mathscr J=0\), this result will be applied with
\(\mathscr J=\mathscr J\mathscr O_Y\) (respectively
\(\mathscr I=\mathscr I\mathscr O_Y\)).}
\textit{Let \(Y\) and \(X\) be two \(S\)-schemes, let \(\mathscr J\)
be a quasi-coherent square-zero ideal of \(\mathscr O_Y\), let
\(Y_{\mathscr J}\) be the closed subscheme of \(Y\) defined by
\(\mathscr J\), and let \(g_{\mathscr J}:Y_{\mathscr J}\to X\) be an
\(S\)-morphism.}
\begin{enumerate}
\item[\SGAInlineRefTarget{sga3:iii:theorem:0-1-8:item:a}(a)] \textit{The set \(P(g_{\mathscr J})\) of \(S\)-morphisms
  \(g:Y\to X\) extending \(g_{\mathscr J}\) is either empty or
  principal homogeneous under the abelian group}
  \[
   \operatorname{Hom}_{\mathscr O_{Y_{\mathscr J}}}
   \bigl(g_{\mathscr J}^*(\Omega^1_{X/S}),\mathscr J\bigr).
  \]
\item[\SGAInlineRefTarget{sga3:iii:theorem:0-1-8:item:b}(b)] \textit{Suppose that \(i:Y_0\hookrightarrow Y_{\mathscr J}\)
  is the closed immersion defined by a quasi-coherent ideal
  \(\mathscr I\supset\mathscr J\) such that
  \(\mathscr I\mathscr J=0\), and set
  \(g_0=g_{\mathscr J}\circ i\). Then the preceding abelian group is
  isomorphic to}
  \[
   \operatorname{Hom}_{\mathscr O_{Y_0}}
   \bigl(g_0^*(\Omega^1_{X/S}),\mathscr J\bigr).
  \]
\end{enumerate}

\noindent\textit{Proof.}
Part~\SGARefLink{sga3:iii:theorem:0-1-8:item:b}{(b)} follows at once from part~\SGARefLink{sga3:iii:theorem:0-1-8:item:a}{(a)}. Indeed, since \(\mathscr J\) is
annihilated by \(\mathscr I\), it may be regarded as an
\(\mathscr O_{Y_0}\)-module, and adjunction gives
\[
 \operatorname{Hom}_{\mathscr O_{Y_{\mathscr J}}}
  \bigl(g_{\mathscr J}^*(\Omega^1_{X/S}),\mathscr J\bigr)
 =
 \operatorname{Hom}_{\mathscr O_{Y_0}}
  \bigl(i^*g_{\mathscr J}^*(\Omega^1_{X/S}),\mathscr J\bigr).
\]

To prove part~\SGARefLink{sga3:iii:theorem:0-1-8:item:a}{(a)}, we may suppose that
\(P(g_{\mathscr J})\ne\emptyset\), that is, that an \(S\)-morphism
\(g:Y\to X\) extending \(g_{\mathscr J}\) exists. Write
\(j:Y_{\mathscr J}\hookrightarrow Y\) for the immersion. Then
\(P(g_{\mathscr J})\) is the set of \(S\)-morphisms \(g':Y\to X\)
such that \(g'\circ j=g_{\mathscr J}\). Giving such a \(g'\) is
equivalent to giving an \(S\)-morphism
\[
 h:Y\longrightarrow X\times_S X
\]
such that \(\operatorname{pr}_1\circ h=g\) and
\(h_{\mathscr J}=\delta\circ g_{\mathscr J}\), where
\(h_{\mathscr J}=h\circ j\) and
\(\delta:X\hookrightarrow X\times_S X\) is the diagonal immersion.

\SGAThreeNativeDiagram[0.56\textwidth]
  {native_diagrams/exp3/diagram_010_exp3_p008_theorem018_factorization.tex}
  {The morphism \(h\) determined by two extensions of \(g_{\mathscr J}\).}
  {fig:exp3-theorem-018-factorization}
  {Polo--Gille Exp3-14oct24.pdf, component-local p. 8; combined-reader p. 118}

Since \(h_{\mathscr J}\) factors through \(\delta\), and since \(Y\)
lies in the first infinitesimal neighborhood of the immersion
\(j:Y_{\mathscr J}\to Y\) because \(\mathscr J^2=0\), functoriality
(EGA~IV\textsubscript{4}, 16.2.2(i)) shows that the desired morphisms
\(h\) factor uniquely through \(\Delta^{(1)}_{X/S}\); see 0.1.7. Put
\[
 Y'=\Delta^{(1)}_{X/S}\times_XY,
 \qquad
 Y'_{\mathscr J}=Y'\times_YY_{\mathscr J}
 =\Delta^{(1)}_{X/S}\times_XY_{\mathscr J}.
\]
The desired morphisms \(h\) are then in bijection with the sections
\(u\) of \(Y'\to Y\) extending the section
\(u_{\mathscr J}=(\delta\circ g_{\mathscr J},\operatorname{id})\) of
\(Y'_{\mathscr J}\to Y_{\mathscr J}\). On the other hand, \(Y'\)
(respectively \(Y'_{\mathscr J}\)) is affine over \(Y\)
(respectively \(Y_{\mathscr J}\)) and corresponds to the
quasi-coherent algebra
\[
 \mathscr A=\mathscr O_Y\oplus g^*(\Omega^1_{X/S}),
 \qquad\text{respectively}\qquad
 \mathscr A_{\mathscr J}
 =\mathscr A\otimes_{\mathscr O_Y}\mathscr O_{Y_{\mathscr J}}
 =\mathscr O_{Y_{\mathscr J}}
   \oplus g_{\mathscr J}^*(\Omega^1_{X/S}).
\]
Let \(\varepsilon:\mathscr A\to\mathscr O_Y\) be the canonical
augmentation, that is, the \(\mathscr O_Y\)-algebra morphism that
vanishes on \(g^*(\Omega^1_{X/S})\), and define
\(\varepsilon_{\mathscr J}:\mathscr A_{\mathscr J}
 \to\mathscr O_{Y_{\mathscr J}}\) similarly. Then
\[
 \Gamma(Y'/Y)
 \simeq\operatorname{Hom}_{\mathscr O_Y\text{-}\mathrm{alg}}
       (\mathscr A,\mathscr O_Y),
 \qquad
 \Gamma(Y'_{\mathscr J}/Y_{\mathscr J})
 \simeq
 \operatorname{Hom}_{\mathscr O_{Y_{\mathscr J}}\text{-}\mathrm{alg}}
       (\mathscr A_{\mathscr J},\mathscr O_{Y_{\mathscr J}}).
\]

Under these isomorphisms, the section
\(u=(\delta\circ g,\operatorname{id})\) (respectively
\(u_{\mathscr J}\)) corresponds to \(\varepsilon\) (respectively
\(\varepsilon_{\mathscr J}\)). Consequently,
\(P(g_{\mathscr J})\) is in bijection with the algebra morphisms
\(\mathscr A\to\mathscr O_Y\) that reduce to
\(\varepsilon_{\mathscr J}\); under this bijection, \(g\) corresponds
to \(\varepsilon\).

Put \(M=g^*(\Omega^1_{X/S})\). Then
\(\operatorname{Hom}_{\mathscr O_Y\text{-}\mathrm{alg}}
 (\mathscr A,\mathscr O_Y)\) identifies with the set of
\(\mathscr O_Y\)-morphisms \(\psi:M\to\mathscr O_Y\) whose image is a
square-zero ideal. We are interested in those morphisms that induce the
zero morphism
\(M\to\mathscr O_{Y_{\mathscr J}}=\mathscr O_Y/\mathscr J\), that is,
those that carry \(M\) into \(\mathscr J\). Conversely, since
\(\mathscr J^2=0\), every \(\mathscr O_Y\)-morphism
\(\phi:M\to\mathscr J\) comes from a unique algebra morphism
\(\mathscr A\to\mathscr O_Y\) reducing to
\(\varepsilon_{\mathscr J}\). Finally,
\[
 \operatorname{Hom}_{\mathscr O_Y}
   \bigl(g^*(\Omega^1_{X/S}),\mathscr J\bigr)
 =
 \operatorname{Hom}_{\mathscr O_{Y_{\mathscr J}}}
   \bigl(g_{\mathscr J}^*(\Omega^1_{X/S}),\mathscr J\bigr)
\]
because \(\mathscr J^2=0\), as in the proof of part~\SGARefLink{sga3:iii:theorem:0-1-8:item:b}{(b)}. Thus one
obtains a bijection
\[
 P(g_{\mathscr J})\simeq
 \operatorname{Hom}_{\mathscr O_{Y_{\mathscr J}}}
 \bigl(g_{\mathscr J}^*(\Omega^1_{X/S}),\mathscr J\bigr)
\]
under which \(g\) corresponds to the zero morphism.

For every
\[
 m\in
 \operatorname{Hom}_{\mathscr O_{Y_{\mathscr J}}}
 \bigl(g_{\mathscr J}^*(\Omega^1_{X/S}),\mathscr J\bigr),
\]
write \(m\cdot g\) for the element of \(P(g_{\mathscr J})\) associated
with \(g\) and \(m\) by the preceding bijection. We have already seen
that \(0\cdot g=g\); it remains to prove
\begin{equation}
 m'\cdot(m\cdot g)=(m+m')\cdot g.
 \tag{0.1.8\,(*)}\label{sga3:iii:theorem:0-1-8:equation:0-1-8-star}
\end{equation}
This may be checked locally.\footnote{Editor's note: The preceding
argument is reproduced almost word for word from SGA~1, III~5.1; what
follows has been expanded.} Indeed, the two preceding morphisms
\(Y\to X\) induce the same continuous map as \(g\) on the underlying
topological spaces. It is therefore enough to check that, for every
affine open \(U=\operatorname{Spec}(A)\) of \(X\) lying over an affine
open \(\operatorname{Spec}(\Lambda)\) of \(S\), and every affine open
\(V=\operatorname{Spec}(B)\) of \(g^{-1}(U)\), they induce the same
morphism of \(\Lambda\)-algebras \(A\to B\).

Set \(J=\Gamma(V,\mathscr J)\), and let
\(\varphi,\psi,\eta:A\to B\) be the morphisms induced respectively by
\(g\), \(m\cdot g\), and \(m'\cdot(m\cdot g)\). They agree modulo
\(J\). There are unique expressions
\[
 \psi=\varphi+D,
 \qquad
 \eta=\psi+D',
\]
where \(D\) belongs to
\[
 \operatorname{Der}_{\varphi}(A,J)
 =\{\delta\in\operatorname{Hom}_{\Lambda}(A,J)
   \mid \delta(ab)=\varphi(a)\delta(b)+\varphi(b)\delta(a)\},
\]
and \(D'\) belongs to \(\operatorname{Der}_{\psi}(A,J)\). But
\(\operatorname{Der}_{\varphi}(A,J)
 =\operatorname{Der}_{\psi}(A,J)\), since \(J^2=0\), and both identify
with
\[
 \operatorname{Hom}_{B/J}
 \bigl(\Omega^1_{A/\Lambda}\otimes_AB/J,J\bigr).
\]
Under this identification, \(D\) corresponds to \(m\) and \(D'\) to
\(m'\). Hence \(\eta=\varphi+D+D'\), while \(D+D'\) corresponds to
\(m+m'\), proving (\SGARefLink{sga3:iii:theorem:0-1-8}{0.1.8}\,(*)).

\medskip
\noindent\SGAInlineRefTarget{sga3:iii:corollary:0-1-9}\textbf{Corollary 0.1.9.}\footnote{Editor's note: This
corollary has been added to SGA~1, III~5.1; it proves part~(i) of
\SGARefLink{sga3:iii:proposition:0-2}{Proposition~0.2} below.}
\textit{Let \(X\) be an \(S\)-scheme, and retain the notation of
\SGARefLink{sga3:iii:definition:0-1-5}{0.1.5}. Then \(X\) carries a left action of the abelian
\(X^+\)-group \(L_X\), making \(X\) formally principal homogeneous
under \(L_X\) over \(X^+\). In other words, there is an isomorphism of
\(X^+\)-functors}
\[
 L_X\times_{X^+}X\xrightarrow{\sim}X\times_{X^+}X,
 \qquad
 (m,x)\longmapsto(x,m\cdot x).
\]

\noindent\textit{Proof.}
Let \(i_0:X_0\hookrightarrow X\) be the immersion. Since
\(X_0=X\times_SS_0\), one first has
\[
 i_0^*(\Omega^1_{X/S})\simeq\Omega^1_{X_0/S_0}
\]
(EGA~IV, 16.4.5).

Let \(Y\) be an \(X^+\)-scheme, given by an \(S\)-morphism
\(g_{\mathscr J}:Y_{\mathscr J}\to X\), and let
\(g_0:Y_0\to X_0\) be obtained by base change. By \SGARefLink{sga3:iii:theorem:0-1-8}{0.1.8}, if
\(\operatorname{Hom}_{X^+}(Y,X)\) is nonempty, it is principal
homogeneous under
\[
 \operatorname{Hom}_{\mathscr O_{Y_0}}
 \bigl(g_0^*i_0^*(\Omega^1_{X/S}),\mathscr J\mathscr O_Y\bigr),
\]
which identifies with
\[
 \operatorname{Hom}_{\mathscr O_{Y_0}}
 \bigl(g_0^*(\Omega^1_{X_0/S_0}),\mathscr J\mathscr O_Y\bigr)
 =L_X(Y).
\]
Thus there is a bijection
\[
 L_X(Y)\times\operatorname{Hom}_{X^+}(Y,X)
 \xrightarrow{\sim}
 \operatorname{Hom}_{X^+}(Y,X\times_{X^+}X),
 \qquad
 (m,g)\longmapsto(g,m\cdot g).
\]
We show that this is functorial in \(Y\). Let \(f:Z\to Y\) be a
morphism of \(S\)-schemes. It is enough to prove that the following
diagram is commutative.

\SGAThreeNativeDiagram[0.88\textwidth]
  {native_diagrams/exp3/diagram_011_exp3_p010_corollary019_functoriality.tex}
  {Functoriality in \(Y\) of the principal-homogeneous bijection.}
  {fig:exp3-corollary-019-functoriality}
  {Polo--Gille Exp3-14oct24.pdf, component-local p. 10; combined-reader p. 120}

If \(\operatorname{Hom}_{X^+}(Y,X)=\emptyset\), there is nothing to
prove. It is therefore enough to show that, for every \(S\)-morphism
\(g:Y\to X\) extending \(g_{\mathscr J}\) and every \(m\in L_X(Y)\),
one has
\begin{equation}
 (m\cdot g)\circ f=L_X(f)(m)\cdot(g\circ f).
 \tag{0.1.9\,(*)}\label{sga3:iii:corollary:0-1-9:equation:0-1-9-star}
\end{equation}
These two \(S\)-morphisms \(Z\to X\) agree on \(Z_{\mathscr J}\) with
\(g_{\mathscr J}\circ f_{\mathscr J}\); in particular, they induce
the same continuous map as \(g\circ f\) on the underlying topological
spaces. It therefore suffices to check the equality on local rings.
Let \(z\in Z\), put \(y=f(z)\), \(x=g(y)\), and let \(A,B,C\) denote
respectively the local rings
\(\mathscr O_{X,x},\mathscr O_{Y,y},\mathscr O_{Z,z}\). Let \(s\) be
the image of \(x\) in \(S\), put \(\Lambda=\mathscr O_{S,s}\), let
\(J\) and \(I\) be the ideals of \(\Lambda\) corresponding to
\(\mathscr J\) and \(\mathscr I\), and let
\(\varphi,\psi:A\to B\) and \(\theta:B\to C\) be the
\(\Lambda\)-algebra morphisms induced by \(g\), \(m\cdot g\), and
\(f\), respectively. Then \(m\) induces an element \(D\) of
\[
 \operatorname{Hom}_{B/IB}
 \bigl(\Omega^1_{A/\Lambda}\otimes_AB/IB,JB\bigr)
 =\operatorname{Der}_{\Lambda}(A,JB),
\]
and \(\psi=\varphi+D\). Thus \((m\cdot g)\circ f\) corresponds to
\[
 \theta\circ\psi=\theta\circ\varphi+\theta\circ D.
\]
By \SGARefLink{sga3:iii:definition:0-1-5}{0.1.5}, \(\theta\circ D\) is the image of \(L_X(f)(m)\) in
\[
 \operatorname{Hom}_{C/IC}
 \bigl(\Omega^1_{A/\Lambda}\otimes_AC/IC,JC\bigr)
 =\operatorname{Der}_{\Lambda}(A,JC).
\]
Consequently, \(\theta\circ\varphi+\theta\circ D\) is the image of
\(L_X(f)(m)\cdot(g\circ f)\). This proves (\SGARefLink{sga3:iii:corollary:0-1-9}{0.1.9}\,(*)).

\medskip
\noindent\SGAInlineRefTarget{sga3:iii:corollary:0-1-10}\textbf{Corollary 0.1.10.}\footnote{Editor's note: The
functoriality of \(L_X\) in \(X\) has been made explicit, in particular
in part~\SGARefLink{sga3:iii:corollary:0-1-10:item:b}{(b)} below.}
\begin{enumerate}
\item[\SGAInlineRefTarget{sga3:iii:corollary:0-1-10:item:a}(a)] \textit{The functor \(L_X\) depends functorially on
  \(X\). For every \(S\)-morphism \(f:X\to W\), there is an
  \(S\)-morphism \(L_f:L_X\to L_W\) that is a morphism of abelian
  groups over \(f^+\); that is, the following diagram is commutative,}

  \SGAThreeNativeDiagram[0.46\textwidth]
    {native_diagrams/exp3/diagram_012_exp3_p011_corollary0110_lf_square.tex}
    {The morphism \(L_f\) over \(f^+\).}
    {fig:exp3-corollary-0110-lf-square}
    {Polo--Gille Exp3-14oct24.pdf, component-local p. 11; combined-reader p. 121}

  \textit{and, for every \(Y\to X^+\), the map}
  \[
   L_f(Y):\operatorname{Hom}_{X^+}(Y,L_X)
   \longrightarrow\operatorname{Hom}_{W^+}(Y,L_W),
  \]
  \textit{where \(Y\) lies over \(W^+\) through \(f^+\), is a
  homomorphism of abelian groups.}

\item[\SGAInlineRefTarget{sga3:iii:corollary:0-1-10:item:b}(b)] \textit{Moreover, the following diagram is commutative.}

  \SGAThreeNativeDiagram[0.66\textwidth]
    {native_diagrams/exp3/diagram_013_exp3_p011_corollary0110_action_square.tex}
    {Compatibility of \(L_f\) with the formally principal homogeneous actions.}
    {fig:exp3-corollary-0110-action-square}
    {Polo--Gille Exp3-14oct24.pdf, component-local p. 11; combined-reader p. 121}
\end{enumerate}

\noindent\textit{Proof.}
For part~\SGARefLink{sga3:iii:corollary:0-1-10:item:a}{(a)}, \(L_f\) is induced by the morphism of
\(\mathscr O_{X_0}\)-modules
\[
 f_{X_0/W_0/S_0}:f_0^*(\Omega^1_{W_0/S_0})
 \longrightarrow\Omega^1_{X_0/S_0}
\]
(cf. \SGARefLink{sga3:iii:paragraph:0-1-7:item:b}{0.1.7(b)}). For every \(X^+\)-scheme \(Y\), given by an
\(S\)-morphism \(g_{\mathscr J}:Y_{\mathscr J}\to X\), one has the
following commutative diagram, functorial in \(Y\).

\SGAThreeNativeDiagram[0.96\textwidth]
  {native_diagrams/exp3/diagram_014_exp3_p011_corollary0110_hom_square.tex}
  {The homomorphism \(L_f(Y)\) on the differential-module descriptions.}
  {fig:exp3-corollary-0110-hom-square}
  {Polo--Gille Exp3-14oct24.pdf, component-local p. 11; combined-reader p. 121}

Here \(L_f(Y)\) is the map
\[
 \psi\longmapsto
 \psi\circ g_0^*(f_{X_0/W_0/S_0}),
\]
which is indeed a homomorphism of abelian groups.\footnote{Editor's
note: Observe that \(L_f\) depends only on
\(f_0:X_0\to W_0\).}

We now prove part~\SGARefLink{sga3:iii:corollary:0-1-10:item:b}{(b)}. Let \(Y\) be an \(X^+\)-scheme. If
\(\operatorname{Hom}_{X^+}(Y,X)=\emptyset\), there is nothing to
prove. Thus let \(g\in\operatorname{Hom}_{X^+}(Y,X)\). We must show
that, for every \(m\in L_X(Y)\),
\begin{equation}
 f\circ(m\cdot g)=L_f(Y)(m)\cdot(f\circ g).
 \tag{0.1.10\,(*)}\label{sga3:iii:corollary:0-1-10:equation:0-1-10-star}
\end{equation}
With \(g\) fixed,
\(\operatorname{Hom}_X(Y,X\times_{X^+}X)\) is a subset of
\(\operatorname{Hom}_X(Y,\Delta^{(1)}_{X/S})\), and
\[
\begin{aligned}
 L_X(Y)
 &=\operatorname{Hom}_{\mathscr O_{Y_0}}
   \bigl(g_0^*(\Omega^1_{X_0/S_0}),\mathscr J\mathscr O_Y\bigr)\\
 &=\operatorname{Hom}_{\mathscr O_Y}
   \bigl(g^*(\Omega^1_{X/S}),\mathscr J\mathscr O_Y\bigr),
\end{aligned}
\]
which is a subset of \(\widetilde L_X(Y)\); see \SGARefLink{sga3:iii:paragraph:0-1-7}{0.1.7}. Moreover,
\(L_f(Y)\) is the restriction of \(\widetilde L_f(Y)\) to
\(L_X(Y)\). Finally, the bijection
\[
 L_X(Y)\xrightarrow{\sim}
 \operatorname{Hom}_X(Y,X\times_{X^+}X),
 \qquad
 m\longmapsto(g,m\cdot g),
\]
is the inverse of the restriction to
\(L_X(Y)\subset\widetilde L_X(Y)\) of the bijection
\[
 \operatorname{Hom}_X(Y,\Delta^{(1)}_{X/S})
 \xrightarrow{\sim}\{g\}\times\widetilde L_X(Y)
\]
considered in \SGARefLink{sga3:iii:remark:0-1-7-1}{0.1.7.1}. Consequently, equality (\SGARefLink{sga3:iii:corollary:0-1-10}{0.1.10}\,(*)) follows
from (\SGARefLink{sga3:iii:paragraph:0-1-7}{0.1.7}\,(*)). Indeed, let \(g'\) be the \(X\)-morphism
\(Y\to\Delta^{(1)}_{X/S}\) defined by \((g,m\cdot g)\). The element
of \(\widetilde L_W(Y)\) corresponding to
\((f\circ g,f\circ g')\) is
\(\widetilde L_f(Y)(m)=L_f(Y)(m)\). Thus
\[
 L_f(m)\cdot(f\circ g)=f\circ(m\cdot g).
\]

